\section{实验分析、局限性与结论}

\subsection{实验完整性}
本文所有最终模型选择均在验证集完成。测试集只在最终方案固定后评估一次；附件 3 和附件 4 为无标签专项测试集，只进行推理。三个 seed 为 42、123、777，最终分类采用三 seed 概率平均后的 Argmax。

\subsection{主要结果}
最终加权交叉熵 DRM-Net 在测试集上的 Accuracy 为 66.30\%，Macro-F1 为 62.54\%，MAE 为 0.6346，Pearson $r$ 为 0.6630。普通交叉熵和 Focal Loss 的 Accuracy 较高，但 Macro-F1 较低，因此最终方案优先保留类别平衡更好的加权交叉熵结果。

\subsection{结构消融结论}
去除缺失增强后 Accuracy 降至 64.65\%，Macro-F1 降至 60.83\%；去除缺失 mask 后 Accuracy 降至 65.75\%，Macro-F1 降至 60.93\%。这说明显式缺失建模是当前模型鲁棒性的主要来源之一。跨模态注意力和动态门控的作用通过独立结构消融报告，不将单一指标解释为绝对优越性。

\subsection{误差与解释边界}
验证集错误主要集中在弱情感混淆、中性边界漂移和严重极性倒置。由于情感标签具有主观性，模型在中性附近存在不可完全消除的歧义。模态遮挡敏感性反映的是模型输出对输入扰动的响应，受模态交互、遮挡方式和模型参数影响，不能作为严格因果证明。

\subsection{局限性}
第一，当前实验结论只覆盖已经实际运行的缺失类型、位置和比例，不能外推到未测试的极端丢包条件。第二，附件 3 和附件 4 没有标签，专项结果不能用于性能验证。第三，当前解释方法是遮挡重算，不是积分梯度或 Attention Rollout。第四，统一长度处理是工程化时间索引，不等价于动态时间规整，因此不声称获得 DTW 的理论保证。

\subsection{结论}
本文完成了从三模态特征提取、时间统一、局部缺失鲁棒训练，到无标签专项预测和遮挡敏感性解释的完整流程。代码、CSV、JSON 和图表均保留了样本级追溯信息，最终结论以当前冻结的实验结果为准。

